% Options for packages loaded elsewhere
\PassOptionsToPackage{unicode}{hyperref}
\PassOptionsToPackage{hyphens}{url}
%
\documentclass[
]{article}
\usepackage{amsmath,amssymb}
\usepackage{iftex}
\ifPDFTeX
  \usepackage[T5]{fontenc}
  \usepackage[utf8]{inputenc}
  \usepackage{textcomp} % provide euro and other symbols
\else % if luatex or xetex
  \usepackage{unicode-math} % this also loads fontspec
  \defaultfontfeatures{Scale=MatchLowercase}
  \defaultfontfeatures[\rmfamily]{Ligatures=TeX,Scale=1}
\fi
\usepackage{lmodern}
\ifPDFTeX\else
  % xetex/luatex font selection
\fi
% Use upquote if available, for straight quotes in verbatim environments
\IfFileExists{upquote.sty}{\usepackage{upquote}}{}
\IfFileExists{microtype.sty}{% use microtype if available
  \usepackage[]{microtype}
  \UseMicrotypeSet[protrusion]{basicmath} % disable protrusion for tt fonts
}{}
\makeatletter
\@ifundefined{KOMAClassName}{% if non-KOMA class
  \IfFileExists{parskip.sty}{%
    \usepackage{parskip}
  }{% else
    \setlength{\parindent}{0pt}
    \setlength{\parskip}{6pt plus 2pt minus 1pt}}
}{% if KOMA class
  \KOMAoptions{parskip=half}}
\makeatother
\usepackage{xcolor}
\usepackage{graphicx}
\makeatletter
\def\maxwidth{\ifdim\Gin@nat@width>\linewidth\linewidth\else\Gin@nat@width\fi}
\def\maxheight{\ifdim\Gin@nat@height>\textheight\textheight\else\Gin@nat@height\fi}
\makeatother
% Scale images if necessary, so that they will not overflow the page
% margins by default, and it is still possible to overwrite the defaults
% using explicit options in \includegraphics[width, height, ...]{}
\setkeys{Gin}{width=\maxwidth,height=\maxheight,keepaspectratio}
% Set default figure placement to htbp
\makeatletter
\def\fps@figure{htbp}
\makeatother
\ifLuaTeX
  \usepackage{luacolor}
  \usepackage[soul]{lua-ul}
\else
  \usepackage{soul}
\fi
\setlength{\emergencystretch}{3em} % prevent overfull lines
\providecommand{\tightlist}{%
  \setlength{\itemsep}{0pt}\setlength{\parskip}{0pt}}
\setcounter{secnumdepth}{-\maxdimen} % remove section numbering
\ifLuaTeX
  \usepackage{selnolig}  % disable illegal ligatures
\fi
\usepackage{bookmark}
\IfFileExists{xurl.sty}{\usepackage{xurl}}{} % add URL line breaks if available
\urlstyle{same}
\hypersetup{
  hidelinks,
  pdfcreator={LaTeX via pandoc}}

\author{}
\date{}

\usepackage[vietnamese]{babel}
\usepackage{ulem}
\let\ul\uline
\begin{document}
\ul{Buổi 04 -- Ngày 09-08-2026 -- môn Đại số tuyến tính -- lớp
MA003.F32.LT.CNTT}

\ul{TH4}: $A$ là ma
trận vuông cấp $n$,
với $n \ge 3$

\ul{Ký hiệu}:

$A(i|j) =$ma trận
$A$ xóa bỏ đi dòng
$(i)$ và cột
$(j)$

$c_{ij}^A = (-1)^{i+j} \det[A(i|j)]$, với
$1 \le i, j \le n$

\ul{Ví dụ mẫu}: Cho
$$A = \begin{pmatrix} -5 & 1 & 4 \\ -7 & -2 & 6 \\ 9 & 10 & -3 \end{pmatrix}$$. Tìm
$$\begin{cases} A(2|3) =? \Rightarrow c_{23}^A =? \\ A(3|1) =? \Rightarrow c_{31}^A =? \\ A(1|2) =? \Rightarrow c_{12}^A =? \end{cases}$$

\ul{Giải}:

Ta có:

\begin{quote}
$$A(2|3) = \begin{pmatrix} -5 & 1 \\ 9 & 10 \end{pmatrix} \Rightarrow c_{23}^A = (-1)^{2+3} \det[A(2|3)] = - \begin{vmatrix} -5 & 1 \\ 9 & 10 \end{vmatrix} = -(-50-9) = 59$$
\end{quote}

$$A(3|1) = \begin{pmatrix} 1 & 4 \\ -2 & 6 \end{pmatrix} \Rightarrow c_{31}^A = (-1)^{3+1} \begin{vmatrix} 1 & 4 \\ -2 & 6 \end{vmatrix} = 6+8=14$$

$$A(1|2) = \begin{pmatrix} -7 & 6 \\ 9 & -3 \end{pmatrix} \Rightarrow c_{12}^A = (-1)^{1+2} \begin{vmatrix} -7 & 6 \\ 9 & -3 \end{vmatrix} = -(21-54)=33$$

\ul{Ví dụ 1}: Cho
$$A = \begin{pmatrix} 4 & -3 & 10 \\ 6 & -1 & 7 \\ 9 & 5 & -3 \end{pmatrix}$$. Tìm
$$\begin{cases} A(2|1) =? \Rightarrow c_{21}^A =? \\ A(3|1) =? \Rightarrow c_{31}^A =? \\ A(1|1) =? \Rightarrow c_{11}^A =? \end{cases}$$

\ul{Ví dụ 2}: Cho
$$B = \begin{pmatrix} -8 & 3 & -7 \\ 6 & 10 & -5 \\ 9 & -2 & 1 \end{pmatrix}$$. Tìm
$$\begin{cases} B(2|3) =? \Rightarrow c_{23}^B =? \\ B(3|3) =? \Rightarrow c_{33}^B =? \\ B(1|3) =? \Rightarrow c_{13}^B =? \end{cases}$$

\ul{Ví dụ 3}: Cho
$$C = \begin{pmatrix} -4 & 2 & -9 \\ 10 & 8 & -3 \\ 1 & 6 & -5 \end{pmatrix}$$. Tìm
$$\begin{cases} C(2|1) =? \Rightarrow c_{21}^C =? \\ C(2|2) =? \Rightarrow c_{22}^C =? \\ C(3|2) =? \Rightarrow c_{32}^C =? \end{cases}$$

\ul{Ví dụ 4}: Cho
$$D = \begin{pmatrix} -3 & 1 & 8 & -6 \\ 9 & 0 & -2 & 5 \\ -3 & 7 & -1 & 8 \\ -10 & 6 & 4 & -2 \end{pmatrix}$$. Tìm
$$\begin{cases} D(1|4) =? \Rightarrow c_{14}^D =? \\ D(2|3) =? \Rightarrow c_{23}^D =? \end{cases}$$

\ul{Ví dụ 5}: Cho
$$E = \begin{pmatrix} -2 & 5 & 10 & 0 \\ 6 & -7 & 1 & -3 \\ 4 & 9 & -8 & 2 \\ 5 & -1 & 3 & 7 \end{pmatrix}$$. Tìm
$$\begin{cases} E(3|4) =? \Rightarrow c_{34}^E =? \\ E(3|2) =? \Rightarrow c_{42}^E =? \end{cases}$$

* Tính $\det(A) = |A|$

Ta có thể tính định thức của
$A$ dựa theo bất kỳ
dòng nào hoặc cột nào của
$A$. Cụ thể như
sau:

Xét $A = (a_{ij})_{1 \le i, j \le n}$ thì

$\det(A) = |A| = \xrightarrow{\text{dòng 1}} a_{11} c_{11}^A + a_{12} c_{12}^A + \dots + a_{1n} c_{1n}^A$

$= \xrightarrow{\text{dòng 2}} a_{21} c_{21}^A + a_{22} c_{22}^A + \dots + a_{2n} c_{2n}^A$

\ldots{}

$= \xrightarrow{\text{cột 1}} a_{11} c_{11}^A + a_{21} c_{21}^A + \dots + a_{n1} c_{n1}^A$

$= \xrightarrow{\text{cột 3}} a_{13} c_{13}^A + a_{23} c_{23}^A + \dots + a_{n3} c_{n3}^A$

\ul{Ví dụ mẫu}: Cho
$$A = \begin{pmatrix} -3 & 5 & -2 \\ 10 & 6 & -1 \\ 4 & -5 & 8 \end{pmatrix}$$

Ta có: $\det(A) = |A|$ (dòng 3)
=

$$4c_{31}^A + (-5)c_{32}^A + 8c_{33}^A = 4(-1)^{3+1} \begin{vmatrix} 5 & -2 \\ 6 & -1 \end{vmatrix} + (-5)(-1)^{3+2} \begin{vmatrix} -3 & -2 \\ 10 & -1 \end{vmatrix} + 8(-1)^{3+3} \begin{vmatrix} -3 & 5 \\ 10 & 6 \end{vmatrix}$$

$= 4(-5+12) + 5(3+20) + 8(-18-50) = -401$

\emph{\ul{Lưu ý}}:

+ Ta thường chọn dòng hoặc cột nào có nhiều hệ số = 0 nhất để tính
$\det(A) = |A|$.

+ Nếu ta thấy trong
$A$ có 2 dòng tỷ lệ
nhau (hoặc 2 cột tỷ lệ nhau) thì suy ra
$\det(A) = |A| = 0$.

\textbf{* \ul{MỐI QUAN HỆ GIỮA CÁC PHÉP BIẾN ĐỔI SƠ CẤP TRÊN DÒNG VÀ
TRÊN CỘT ĐỐI VỚI ĐỊNH THỨC MA TRẬN VUÔNG}}

Cho $A$ là ma trận
vuông cấp $n$, trên
\emph{F}.

Tương tự như 3 phép biến đổi sơ cấp trên dòng, ta cũng có 3 phép biến
đổi sơ cấp trên cột

\textbf{a/ \emph{\ul{Loại 1}}}

Trên dòng Trên cột

$A \xrightarrow{(i) \rightarrow c(i)} A'$
$A \xrightarrow{(i)' \rightarrow c(i)'} A'$

Thì lúc này $\det(A') = c \det(A)$
(gấp $c$ lần)

\textbf{b/ \emph{\ul{Loại 2}}}:

\begin{quote}
Trên dòng Trên cột
\end{quote}

$A \xrightarrow{(i) \leftrightarrow (j)} A'$
$A \xrightarrow{(i)' \leftrightarrow (j)'} A'$

Thì lúc này $\det(A') = -\det(A)$
(ngược dấu)

\textbf{c/ \emph{\ul{Loại 3}}}:

\begin{quote}
Trên dòng Trên cột
\end{quote}

$A \xrightarrow{(i) \rightarrow (i) + c(j)} A'$
$A \xrightarrow{(i)' \rightarrow (i)' + c(j)'} A'$

Thì lúc này $\det(A') = \det(A)$
(không thay đổi)

Cho nên ta có thể vận dụng tùy ý các phép biến đổi sơ cấp trên dòng hoặc
cột (khuyến khích dùng loại 3) để ta đưa 1 dòng hoặc 1 cột nào đó về
nhiều hệ số = 0 nhất, rồi sau đó tính định thức theo dòng hoặc cột có
nhiều hệ số = 0 này.

\ul{Ví dụ mẫu}:

Giải phương trình:
$$\begin{vmatrix} 3 & 3 & x & x \\ x & 3 & 3 & x \\ x & x & 3 & 3 \\ x & x & x & 3 \end{vmatrix} = 0$$ (*)

Ta tính

$$\begin{vmatrix} 3 & 3 & x & x \\ x & 3 & 3 & x \\ x & x & 3 & 3 \\ x & x & x & 3 \end{vmatrix} \xrightarrow[(3) \rightarrow (3)-(2)]{(4) \rightarrow (4)-(3)} \begin{vmatrix} 3 & 3 & x & x \\ x & 3 & 3 & x \\ 0 & x-3 & 0 & 3-x \\ 0 & 0 & x-3 & 0 \end{vmatrix}$$

Tính theo dòng (4), ta có

$$= 0.c_{41}^{A'} + 0.c_{42}^{A'} + (x-3)c_{43}^{A'} + 0.c_{44}^{A'} \\ = (x-3)(-1)^{4+3} \begin{vmatrix} 3 & 3 & x \\ x & 3 & x \\ 0 & x-3 & 3-x \end{vmatrix}$$

$$\xrightarrow{(2)' \rightarrow (2)' + (3)'} (3-x) \begin{vmatrix} 3 & 3+x & x \\ x & 3+x & x \\ 0 & 0 & 3-x \end{vmatrix}$$

Tính theo dòng (3), ta có

$$= (3-x) [0.c_{31}^{A''} + 0.c_{32}^{A''} + (3-x).c_{33}^{A''}] \\ = (3-x)(3-x)(-1)^{3+3} \begin{vmatrix} 3 & 3+x \\ x & 3+x \end{vmatrix} \\ = (3-x)^2 [3(3+x) - x(3+x)] \\ = (3-x)^2 (9 + 3x - 3x - x^2) \\ = (3-x)^2 (9 - x^2) \\ = (3-x)^2 (3-x)(3+x) \\ = (3-x)^3 (3+x)$$

Nên (*) tương đương

$$(3-x)^3(3+x) = 0 \Leftrightarrow \begin{bmatrix} (3-x)^3 = 0 \\ 3+x = 0 \end{bmatrix} \Leftrightarrow \begin{bmatrix} 3-x = 0 \\ 3+x = 0 \end{bmatrix} \Leftrightarrow \begin{bmatrix} x = 3 \\ x = -3 \end{bmatrix}$$

\ul{Đáp số}: pt này có 2 nghiệm là
$x = 3$ hay
$x = -3$.

\ul{Ví dụ 1}: Tính định thức ma trận sau:

$$A = \begin{pmatrix} 1 & 2 & 3 & 4 & 5 \\ 2 & 3 & 7 & 10 & 13 \\ 3 & 5 & 11 & 16 & 21 \\ 2 & -7 & 7 & 7 & 2 \\ 1 & 4 & 5 & 3 & 10 \end{pmatrix}$$

\ul{Ví dụ 2}: Giải phương trình sau trên
$\mathbb{R}$:
$$\begin{vmatrix} x & 1 & 1 & 1 \\ 1 & x & 1 & 1 \\ 1 & 1 & x & 1 \\ 1 & 1 & 1 & x \end{vmatrix} = 0$$ (với
$x$ là ẩn số cần
tìm)

Ta có pt tương đương

\begin{quote}
$(x-1)(x^3-3x+2) + (x-1)(x^2-2x+1) = 0 \\ \Leftrightarrow (x-1)[x^3-3x+2+x^2-2x+1] = 0 \\ \Leftrightarrow (x-1)(x^3+x^2-5x+3) = 0 \\ \Leftrightarrow (x-1)^3(x+3) = 0$
\end{quote}

\textbf{\ul{Lưu ý}}: Khi \emph{A} là ma trận đường chéo, hoặc là ma trận
tam giác trên, hoặc là ma trận tam giác dưới, thì det(\emph{A}) = tích
các phần tử trên đường chéo chính.

\ul{Ví dụ}: Cho
$$A = \begin{pmatrix} 3 & -2 & 1 & 0 \\ 0 & 1 & -7 & 5 \\ 0 & 0 & 6 & -1 \\ 0 & 0 & 0 & -2 \end{pmatrix}$$$\Rightarrow \det(A) = 3 \cdot 1 \cdot 6 \cdot (-2) = -36$.

\textbf{* \ul{MỐI LIÊN HỆ GIỮA MA TRẬN KHẢ NGHỊCH VÀ ĐỊNH THỨC CỦA MA
TRẬN VUÔNG}}

Cho ma trận $A$
vuông cấp $n$ trên
\emph{F}. Ta sẽ kiểm tra xem
$A$ có khả nghịch
hay không? Và tìm ma trận nghịch đảo
$A^{-1}$ (nếu có) bằng
phương pháp định thức như sau:

\ul{Bước 1}: Ta tìm
$\det(A) = |A| = \dots$

\ul{Bước 2}: Nếu
$\det(A) = |A| = 0$ ta kết luận
$A$ không khả
nghịch, và không tồn tại
$A^{-1}$

Nếu $\det(A) = |A| \neq 0$ thì
$A$ khả nghịch, và
ta chuyển sang Bước 3.

\ul{Bước 3}: Ta tìm
$n^2$ hệ số
$c_{ij}^A = (-1)^{i+j} \det[A(i|j)]$, với
$1 \le i, j \le n$.

\ul{Bước 4}: Ta lập ma trận
$C = (c_{ij}^A)_{1 \le i, j \le n}$

\ul{Bước 5}: Ta viết
$C^T = (c_{ji}^A)_{1 \le j, i \le n}$,

\begin{quote}
ta gọi $C^T$ là ma
trận chuyển vị (transposed matrix) của
$C$

= là ma trận liên hợp (adjoint matrix) của
$A$, và kí hiệu
$C^T = \operatorname{adj}(A) = A^*$

Suy ra $A^{-1} = \frac{1}{\det(A)} \cdot C^T$
\end{quote}

\ul{Ví dụ mẫu 1}: Cho
$$A = \begin{pmatrix} -3 & 2 & 10 \\ 8 & -1 & 7 \\ 6 & 4 & -2 \end{pmatrix}$$. Hỏi
$A$ có khả nghịch
không? Nếu có, tìm
$A^{-1}$ bằng pp định
thức.

\ul{Giải}: Ta có

$\det(A) = |A| = [(-3)(-1)(-2) + 2 \cdot 7 \cdot 6 + 10 \cdot 8 \cdot 4] - [10(-1)6 + 2 \cdot 8(-2) + (-3)7 \cdot 4] = 574 \neq 0$

Nên $A$ khả nghịch
và tiếp theo ta tìm 9 hệ số
$c_{ij}^A = (-1)^{i+j} \det[A(i|j)]$, với
$1 \le i, j \le 3$.

$$c_{11}^A = (-1)^{1+1} \begin{vmatrix} -1 & 7 \\ 4 & -2 \end{vmatrix} = 2 - 28 = -26$$

$$c_{12}^A = (-1)^{1+2} \begin{vmatrix} 8 & 7 \\ 6 & -2 \end{vmatrix} = -(-16 - 42) = 58$$

$$c_{13}^A = (-1)^{1+3} \begin{vmatrix} 8 & -1 \\ 6 & 4 \end{vmatrix} = 32 + 6 = 38$$

$$c_{21}^A = (-1)^{2+1} \begin{vmatrix} 2 & 10 \\ 4 & -2 \end{vmatrix} = -(-4 - 40) = 44$$

$$c_{22}^A = (-1)^{2+2} \begin{vmatrix} -3 & 10 \\ 6 & -2 \end{vmatrix} = 6 - 60 = -54$$

$$c_{23}^A = (-1)^{2+3} \begin{vmatrix} -3 & 2 \\ 6 & 4 \end{vmatrix} = -(-12 - 12) = 24$$

$$c_{31}^A = (-1)^{3+1} \begin{vmatrix} 2 & 10 \\ -1 & 7 \end{vmatrix} = 14 + 10 = 24$$

$$c_{32}^A = (-1)^{3+2} \begin{vmatrix} -3 & 10 \\ 8 & 7 \end{vmatrix} = -(-21 - 80) = 101$$

$$c_{33}^A = (-1)^{3+3} \begin{vmatrix} -3 & 2 \\ 8 & -1 \end{vmatrix} = 3 - 16 = -13$$

Ta lập ma trận $$C = (c_{ij}^A)_{1 \le i,j \le 3} = \begin{pmatrix} -26 & 58 & 38 \\ 44 & -54 & 24 \\ 24 & 101 & -13 \end{pmatrix}$$,
suy ra $$C^T = \operatorname{adj}(A) = \begin{pmatrix} -26 & 44 & 24 \\ 58 & -54 & 101 \\ 38 & 24 & -13 \end{pmatrix} = A^*$$.

Và ta có $$A^{-1} = \frac{1}{\det(A)} C^T = \frac{1}{574} \begin{pmatrix} -26 & 44 & 24 \\ 58 & -54 & 101 \\ 38 & 24 & -13 \end{pmatrix}$$.

\ul{Ví dụ mẫu 2}: Xét tính khả nghịch của ma trận
$$A = \begin{pmatrix} 1 & a & a^2 \\ 1 & b & b^2 \\ 1 & c & c^2 \end{pmatrix}$$, với
$a,b,c \in \mathbb{R}$ là các tham số
thực

\ul{Giải}: Ta có:

$$A = \begin{pmatrix} 1 & a & a^2 \\ 1 & b & b^2 \\ 1 & c & c^2 \end{pmatrix} \xrightarrow[(3) \rightarrow (3)-(1)]{(2) \rightarrow (2)-(1)} \begin{pmatrix} 1 & a & a^2 \\ 0 & b-a & b^2-a^2 \\ 0 & c-a & c^2-a^2 \end{pmatrix}$$

Suy ra $\det(A) = |A|$ (cột 1)
=

$$1.c_{11}^{A'} = 1(-1)^{1+1} \begin{vmatrix} b-a & b^2-a^2 \\ c-a & c^2-a^2 \end{vmatrix} = (b-a)(c^2-a^2) - (b^2-a^2)(c-a) = (b-a)(c-a)(c+a-b-a)$$

$= (b-a)(c-a)(c-b)$

Ma trận $A$ khả
nghịch khi

$\det(A) = |A| \neq 0 \Leftrightarrow (b-a)(c-a)(c-b) \neq 0 \Leftrightarrow a \neq b \neq c$

Ma trận $A$ không
khả nghịch khi

$$\det(A) = |A| = 0 \Leftrightarrow (b-a)(c-a)(c-b) = 0 \Leftrightarrow \begin{bmatrix} a = b \\ b = c \\ c = a \end{bmatrix}$$

Nghĩa là $a=b$ hay
$b=c$ hay
$c=a$.

\emph{\ul{Bài tập tương tự}}:

\ul{Bài 1}: Cho
$$A = \begin{pmatrix} 4 & -3 & 2 \\ 10 & -6 & 5 \\ 7 & 1 & -4 \end{pmatrix}$$. Hỏi
$A$ có khả nghịch
không?

Nếu có, tìm $A^{-1}$
bằng pp định thức.

\ul{Bài 2}: Cho
$$A = \begin{pmatrix} 2 & -7 & -9 \\ 8 & -1 & 3 \\ -4 & 9 & -6 \end{pmatrix}$$. Hỏi
$A$ có khả nghịch
không?

Nếu có, tìm $A^{-1}$
bằng pp định thức.

\ul{Bài 3}: Khi nào ma trận
$A$ khả nghịch?
Khi đó, tìm $A^{-1}$
bằng phương pháp định thức, với

$$A = \begin{pmatrix} 1 & a & a \\ a & 1 & a^2 \\ a^2 & a & 1 \end{pmatrix}$$

\ul{Bài 4}: Khi nào ma trận
$A$ khả nghịch?
Khi đó, tìm $A^{-1}$
bằng phương pháp định thức, với

$$A = \begin{pmatrix} 1 & a & a^2 \\ a & a^2 & 1 \\ a^2 & 1 & a \end{pmatrix}$$

\ul{Bài 5}: Khi nào ma trận
$A$ khả nghịch?
Khi đó, tìm $A^{-1}$
bằng phương pháp định thức, với

$$A = \begin{pmatrix} 1 & 1 & 1 \\ a & b & c \\ bc & ac & ab \end{pmatrix}$$

\textbf{* \ul{GIẢI HỆ PHƯƠNG TRÌNH TUYẾN TÍNH BẰNG PHƯƠNG PHÁP ĐỊNH THỨC
(PP CRAMER)}}

Cho hệ phương trình tuyến tính gồm có
$n$ phương trình
và $n$ ẩn số, trên
\emph{F}, viết dưới dạng tích ma trận
$AX = B$, nghĩa là ban
đầu ta có hệ

$$\begin{cases} a_{11} x_1 + a_{12} x_2 + \dots + a_{1n} x_n = b_1 \\ a_{21} x_1 + a_{22} x_2 + \dots + a_{2n} x_n = b_2 \\ \dots \dots \dots \dots \dots \dots \dots \dots \dots \dots \dots \dots \\ a_{n1} x_1 + a_{n2} x_2 + \dots + a_{nn} x_n = b_n \end{cases}$$ (*)

Ta đặt $A = (a_{ij})_{1 \le i, j \le n}$;
$$X = \begin{pmatrix} x_1 \\ x_2 \\ \vdots \\ x_n \end{pmatrix}$$;
$$B = \begin{pmatrix} b_1 \\ b_2 \\ \vdots \\ b_n \end{pmatrix}$$.

Gọi $\Delta = \det(A) = |A|$

$A_j =$ma trận
$A$ sau khi ta xóa
bỏ cột $(j)$ và thay
bằng cột $B$
($1 \le j \le n$)

$\Delta_j = \det(A_j) = |A_j|$

\ul{TH1}: Nếu $\Delta = \det(A) = |A| \neq 0$
thì hệ pt (*) có bộ nghiệm duy nhất là:

$x_1 = \frac{\Delta_1}{\Delta}; x_2 = \frac{\Delta_2}{\Delta}; \dots; x_n = \frac{\Delta_n}{\Delta}$

\ul{TH2}: Nếu $\Delta = \det(A) = |A| = 0$
và có (ít nhất) một hệ số khác 0 trong các số
$\Delta_1, \Delta_2, \dots, \Delta_n$ thì hệ pt (*)
VÔ NGHIỆM

\ul{TH3}: Nếu $\Delta = \det(A) = |A| = 0 = \Delta_1 = \Delta_2 = \dots = \Delta_n$
thì ta kết luận hệ không có nghiệm duy nhất (nghĩa là hệ pt (*) có thể
vô nghiệm hoặc là vô số nghiệm) $\rightarrow$ muốn biết chính xác thì ta phải giải
hệ bằng pp GAUSS-JORDAN hoặc là pp GAUSS.

\ul{Ví dụ mẫu 2}:

Giải hệ pt bằng pp Cramer:

$$\begin{cases} 2x_1 + 2x_2 + 5x_3 = 21 \\ 2x_1 + 3x_2 + 6x_3 = 26 \\ x_1 - 6x_2 - 9x_3 = -37 \end{cases}$$ (*)

Đặt

$$A = \begin{pmatrix} 2 & 2 & 5 \\ 2 & 3 & 6 \\ 1 & -6 & -9 \end{pmatrix}; X = \begin{pmatrix} x_1 \\ x_2 \\ x_3 \end{pmatrix}; B = \begin{pmatrix} 21 \\ 26 \\ -37 \end{pmatrix}$$

Suy ra

$\Delta = \det(A) = |A| = 2 \cdot 3(-9) + 2 \cdot 6 \cdot 1 + 5 \cdot 2(-6) - [5 \cdot 3 \cdot 1 + 2 \cdot 2(-9) + 2 \cdot 6(-6)] = -9 \neq 0$

Nên hệ (*) có nghiệm duy nhất.

$x_1 = \frac{\Delta_1}{\Delta}; x_2 = \frac{\Delta_2}{\Delta}; x_3 = \frac{\Delta_3}{\Delta}$

Trong đó:

$$A_1 = \begin{pmatrix} 21 & 2 & 5 \\ 26 & 3 & 6 \\ -37 & -6 & -9 \end{pmatrix} \Rightarrow \Delta_1 = \det(A_1) = -12$$

$$A_2 = \begin{pmatrix} 2 & 21 & 5 \\ 2 & 26 & 6 \\ 1 & -37 & -9 \end{pmatrix} \Rightarrow \Delta_2 = \det(A_2) = -20$$

$$A_3 = \begin{pmatrix} 2 & 2 & 21 \\ 2 & 3 & 26 \\ 1 & -6 & -37 \end{pmatrix} \Rightarrow \Delta_3 = \det(A_3) = -25$$

Kết luận: hệ (*) có nghiệm duy nhất là:

$x_1 = \frac{\Delta_1}{\Delta} = \frac{-12}{-9} = \frac{4}{3}$

$x_2 = \frac{\Delta_2}{\Delta} = \frac{-20}{-9} = \frac{20}{9}$

$x_3 = \frac{\Delta_3}{\Delta} = \frac{-25}{-9} = \frac{25}{9}$

\ul{Ví dụ}: Cho hệ pt tuyến tính trên
$\mathbb{R}$ như sau:

$$\begin{cases} 2x_2 + 2x_3 + x_1 = 0 \\ (m-2)x_2 - 2x_1 + (m-5)x_3 = 2 \\ x_2 + (m+1)x_3 + mx_1 = -2 \end{cases}$$

với $m \in \mathbb{R}$ là tham số
thực.

Hãy giải và biện luận hệ pt này bằng pp CRAMER.

\ul{Gợi ý}:

Đặt $$A = \begin{pmatrix} 1 & 2 & 2 \\ -2 & m-2 & m-5 \\ m & 1 & m+1 \end{pmatrix}$$;
$$B = \begin{pmatrix} 0 \\ 2 \\ -2 \end{pmatrix}$$;
$$X = \begin{pmatrix} x_1 \\ x_2 \\ x_3 \end{pmatrix}$$.

Gọi

$$\Delta = \det(A) = \begin{vmatrix} 1 & 2 & 2 \\ -2 & m-2 & m-5 \\ m & 1 & m+1 \end{vmatrix} =$$

Ta có:

$$A_1 = \begin{pmatrix} 0 & 2 & 2 \\ 2 & m-2 & m-5 \\ -2 & 1 & m+1 \end{pmatrix} \Rightarrow \Delta_1 = \det(A_1) = \begin{vmatrix} 0 & 2 & 2 \\ 2 & m-2 & m-5 \\ -2 & 1 & m+1 \end{vmatrix} =$$

Và

$$A_2 = \begin{pmatrix} 1 & 0 & 2 \\ -2 & 2 & m-5 \\ m & -2 & m+1 \end{pmatrix} \Rightarrow \Delta_2 = \det(A_2) = \begin{vmatrix} 1 & 0 & 2 \\ -2 & 2 & m-5 \\ m & -2 & m+1 \end{vmatrix} =$$

Và

$$A_3 = \begin{pmatrix} 1 & 2 & 0 \\ -2 & m-2 & 2 \\ m & 1 & -2 \end{pmatrix} \Rightarrow \Delta_3 = \det(A_3) = \begin{vmatrix} 1 & 2 & 0 \\ -2 & m-2 & 2 \\ m & 1 & -2 \end{vmatrix} =$$

Biện luận:

\ul{TH1}: $\Delta \neq 0 \Leftrightarrow \dots \neq 0 \Leftrightarrow \dots$

Thì hệ pt có nghiệm duy nhất là:

$x_1 = \frac{\Delta_1}{\Delta} = \dots ; x_2 = \frac{\Delta_2}{\Delta} = \dots ; x_3 = \frac{\Delta_3}{\Delta} = \dots$

\ul{TH2}: $$\Delta = 0 \Leftrightarrow \begin{bmatrix} m = \dots \\ m = \dots \end{bmatrix}$$

\ul{TH2.1} $m = \dots$
thay vào các delta tìm giá trị:

\begin{quote}
$$\begin{aligned} \Delta_1 &= \dots ? \\ \Delta_2 &= \dots ? \\ \Delta_3 &= \dots ? \end{aligned}$$
\end{quote}

\ul{TH2.2} $m = \dots$
thay vào các delta tìm giá trị:

\begin{quote}
$$\begin{aligned} \Delta_1 &= \dots ? \\ \Delta_2 &= \dots ? \\ \Delta_3 &= \dots ? \end{aligned}$$
\end{quote}

Kết luận\ldots{}

\end{document}
